\documentclass[12pt,a4paper]{article}

% PAQUETES
\usepackage[utf8]{inputenc}
\usepackage[T1]{fontenc}
\usepackage[spanish]{babel}
\usepackage{geometry}
\usepackage{array}
\usepackage{tabularx}
\usepackage{booktabs}
\usepackage{longtable}
\usepackage{enumitem}
\usepackage{float}
\usepackage{placeins}
\usepackage{xcolor}
\usepackage{hyperref}

\geometry{margin=2.5cm}
\renewcommand{\arraystretch}{1.3}
\setlength{\parindent}{0pt}
\setlength{\parskip}{6pt}

\newcolumntype{Y}{>{\raggedright\arraybackslash}X}

% Marca visible para lo que falta. Antes de entregar no debe quedar ninguna:
% buscar "Pendiente" en el PDF y en este archivo.
\newcommand{\pend}[1]{\textcolor{red}{[Pendiente: #1]}}

\begin{document}

% PORTADA
\begin{titlepage}
    \centering

    {\large \textbf{UNIVERSIDAD TÉCNICA ESTATAL DE QUEVEDO}}\\[0.3cm]
    {\large Facultad de Ciencias de la Computación}\\
    {\large Carrera de Software}\\[1.5cm]

    {\Large \textbf{INGENIERÍA DE REQUERIMIENTOS}}\\[0.5cm]
    {\large ISR-401}\\[1.5cm]

    {\LARGE \textbf{GA-TL-02 -- FORM}}\\[0.9cm]

    \rule{0.85\textwidth}{0.5pt}\\[0.5cm]

    {\large \textbf{Taller de entrevista estructurada y design thinking aplicados al PFC}}\\

    \rule{0.85\textwidth}{0.5pt}\\[1.5cm]

    \begin{flushleft}
        \textbf{Integrantes (solo quienes trabajaron):}
        \begin{itemize}[noitemsep,topsep=2pt]
            \item \pend{apellidos y nombres, integrante 1}
            \item \pend{apellidos y nombres, integrante 2}
            \item \pend{apellidos y nombres, integrante 3}
            \item \pend{apellidos y nombres, integrante 4}
        \end{itemize}
        \vspace{0.2cm}
        \textbf{Nivel:} Cuarto\\[0.4cm]
        \textbf{Paralelo:} A\\[0.4cm]
        \textbf{Docente:} Gleiston Cicerón Guerrero Ulloa, PhD\\[0.4cm]
        \textbf{Modalidad:} Grupal\\[0.4cm]
        \textbf{Fecha límite de entrega:} 21 de noviembre de 2026, 23:59
    \end{flushleft}

    \vfill

    {\large Quevedo -- Ecuador}\\
    {\large 2026}

\end{titlepage}

% 1. CONTEXTO
\section{Contexto del ejercicio}

% Responsable: persona 2
Proyecto de fin de carrera (PFC) usado: \pend{nombre del PFC}.

Parte interesada entrevistada: \pend{cargo o rol de la parte interesada}.

Fecha y lugar del ensayo: \pend{fecha y lugar}.

Roles en el ensayo:
\begin{itemize}[noitemsep]
    \item Entrevistador: \pend{nombre}.
    \item Parte interesada (juego de roles): \pend{nombre}.
    \item Registro y observación: \pend{nombre}.
\end{itemize}

% 2. GUION
\section{Guion de entrevista estructurada}

% Responsable: persona 1
% Entre 12 y 15 preguntas, con apertura, bloques temáticos y cierre.
% Preguntas abiertas, neutrales y simples (una idea por pregunta).

\begingroup\small
\begin{longtable}{>{\raggedright\arraybackslash}p{0.8cm} >{\raggedright\arraybackslash}p{2.8cm} >{\raggedright\arraybackslash}p{6.3cm} >{\raggedright\arraybackslash}p{4.0cm}}
\caption{Guion de entrevista}\\
\toprule
\textbf{N.°} & \textbf{Bloque} & \textbf{Pregunta} & \textbf{Qué se quiere saber} \\
\midrule
\endfirsthead
\toprule
\textbf{N.°} & \textbf{Bloque} & \textbf{Pregunta} & \textbf{Qué se quiere saber} \\
\midrule
\endhead
1 & Apertura & \pend{pregunta} & \pend{objetivo} \\
\midrule
2 & Bloque 1: \pend{tema} & \pend{pregunta} & \pend{objetivo} \\
\midrule
3 & Bloque 1: \pend{tema} & \pend{pregunta} & \pend{objetivo} \\
\midrule
% ... agregar filas hasta tener entre 12 y 15 preguntas ...
12 & Cierre & \pend{pregunta} & \pend{objetivo} \\
\bottomrule
\end{longtable}
\endgroup

% 3. CONSENTIMIENTO Y MANIFIESTO
\section{Consentimiento informado y registro de evidencias}

% Responsable: persona 4
\subsection{Protocolo de consentimiento}

\pend{texto del protocolo, una carilla: propósito, qué se graba, uso, derecho a retirarse}

\subsection{Manifiesto de evidencias}

El audio y el consentimiento firmado no se publican en este repositorio.
Se guardan en una carpeta institucional privada (cuenta \texttt{@uteq.edu.ec})
compartida con el docente. Aquí solo se registra el siguiente manifiesto.

\begin{table}[H]
\centering
\small
\begin{tabularx}{\textwidth}{>{\raggedright\arraybackslash}p{3.2cm} p{2.4cm} p{2.0cm} Y}
\toprule
\textbf{Archivo} & \textbf{Fecha} & \textbf{Duración} & \textbf{Huella SHA-256} \\
\midrule
\pend{nombre del archivo de audio} & \pend{fecha} & \pend{mm:ss} & \pend{huella} \\
\bottomrule
\end{tabularx}
\caption{Manifiesto de evidencias}
\end{table}

% 4. REGISTRO DEL ENSAYO
\section{Registro del ensayo}

% Responsable: persona 3
\pend{resumen de lo dicho en el ensayo, en orden, con las frases clave}

% 5. MAPA DE EMPATÍA
\section{Mapa de empatía}

% Responsable: persona 3
% Debe ser coherente con lo dicho en la entrevista.

\begin{table}[H]
\centering
\small
\begin{tabularx}{\textwidth}{>{\raggedright\arraybackslash}p{3.0cm} Y}
\toprule
\textbf{Dimensión} & \textbf{Contenido} \\
\midrule
Qué dice & \pend{...} \\
\midrule
Qué piensa & \pend{...} \\
\midrule
Qué hace & \pend{...} \\
\midrule
Qué siente & \pend{...} \\
\midrule
Dolores & \pend{...} \\
\midrule
Ganancias esperadas & \pend{...} \\
\bottomrule
\end{tabularx}
\caption{Mapa de empatía de \pend{usuario principal}}
\end{table}

\section{Definición del problema}

% Responsable: persona 3
% Centrada en el usuario, no en la solución.
Usuario: \pend{quién}. Necesidad: \pend{qué necesita}. Hallazgo: \pend{por qué}.

\textbf{¿Cómo podríamos} \pend{...}\textbf{?}

% 6. AJUSTES
\section{Ajustes al guion}

% Responsable: persona 4
\begingroup\small
\begin{longtable}{>{\raggedright\arraybackslash}p{4.8cm} >{\raggedright\arraybackslash}p{3.2cm} >{\raggedright\arraybackslash}p{6.7cm}}
\caption{Preguntas que fallaron y su versión corregida}\\
\toprule
\textbf{Pregunta original} & \textbf{Defecto} & \textbf{Versión corregida} \\
\midrule
\endfirsthead
\toprule
\textbf{Pregunta original} & \textbf{Defecto} & \textbf{Versión corregida} \\
\midrule
\endhead
\pend{pregunta} & \pend{inductiva, doble o cerrada sin necesidad} & \pend{corregida} \\
\bottomrule
\end{longtable}
\endgroup

% 7. DECLARACIÓN DE USO DE IA
\section{Declaración de uso de herramientas de IA}

% Responsable: todos. Si no se usó ninguna herramienta, escribirlo así.
\pend{herramienta, versión, propósito y cómo se verificó lo obtenido; o ``No se usó ninguna herramienta de IA''}

% 8. PARTICIPACIÓN
\section{Participación de los integrantes}

\begin{table}[H]
\centering
\small
\begin{tabularx}{\textwidth}{>{\raggedright\arraybackslash}p{4.5cm} Y}
\toprule
\textbf{Integrante} & \textbf{Parte realizada} \\
\midrule
\pend{nombre} & \pend{parte} \\
\bottomrule
\end{tabularx}
\caption{Reparto del trabajo}
\end{table}

\end{document}
